\chapter{Học sinh hỏi về điểm số và chuyện xin điểm}
\markboth{Học sinh hỏi về điểm số và chuyện xin điểm}{Học sinh hỏi về điểm số và chuyện xin điểm}

\section{Từ câu 101 đến câu 120}
\begin{enumerate}[leftmargin=*]
\item \hsqa{Thầy cho em thêm điểm đi.}{Thầy không bán điểm, thầy chỉ chấm bài.}
\item \hsqa{Em thiếu có nửa điểm thôi.}{Nửa điểm với ranh giới thường rất lạnh lùng.}
\item \hsqa{Thầy thương em cộng chút được không?}{Thương thì nhắc em học, chứ không bẻ thang điểm.}
\item \hsqa{Bài em làm gần đúng mà thầy.}{Gần đúng là chưa đủ để gặp điểm cao.}
\item \hsqa{Sao bài em ít điểm quá?}{Vì đáp án của em cũng ít đúng quá.}
\item \hsqa{Em học mà điểm thấp.}{Em học thật hay em thấy mình ngồi với vở?}
\item \hsqa{Bạn em cũng sai mà được hơn em.}{Bạn em sai chỗ khác, em sai chỗ đắt hơn.}
\item \hsqa{Thầy ơi em bị oan.}{Bài thi ít khi oan bằng lời kể sau giờ.}
\item \hsqa{Em thấy điểm này hơi đau.}{Điểm đau nhất lúc nó đúng thực lực.}
\item \hsqa{Em làm hết mà sao không cao?}{Làm hết không bằng làm đúng.}
\item \hsqa{Em có thể xin chấm nương tay không?}{Có thể xin, còn thầy có nương không là chuyện khác.}
\item \hsqa{Thầy chấm khó quá.}{Thầy chấm theo đáp án, không theo cảm xúc thương trò.}
\item \hsqa{Thầy cho em cơ hội sửa điểm đi.}{Điểm có thể sửa, nhưng trước hết em sửa cách học đã.}
\item \hsqa{Em thiếu mỗi kết luận thôi.}{Mỗi kết luận đôi khi kéo cả bài xuống.}
\item \hsqa{Em nghĩ bài này đáng hơn vậy.}{Bài em nghĩ cao thường không quyết định được điểm em nhận.}
\item \hsqa{Em làm ra ý đó mà thầy.}{Ra ý chưa ra điểm, còn phải ra cách trình bày nữa.}
\item \hsqa{Điểm em thấp vì run.}{Run là lý do thật, nhưng bài vẫn chấm theo chữ em viết.}
\item \hsqa{Thầy coi lại bài em đi.}{Được, miễn em cũng coi lại cách em làm.}
\item \hsqa{Em bị hụt mất câu cuối.}{Câu cuối hay là câu lấy đi nụ cười đó em.}
\item \hsqa{Em đáng được hơn.}{Đáng hay không, bài làm đang trả lời rõ rồi.}
\end{enumerate}

\section{Từ câu 121 đến câu 140}
\begin{enumerate}[leftmargin=*,resume]
\item \hsqa{Sao bài em sạch đẹp mà điểm thấp?}{Sạch đẹp là ưu điểm, nhưng đúng mới là lương thực.}
\item \hsqa{Em tưởng chỗ đó thầy cho điểm.}{Thầy cho điểm cái đúng, không cho điểm niềm tin.}
\item \hsqa{Em làm nhầm một xíu thôi.}{Một xíu đúng chỗ hiểm là bay cả đoạn.}
\item \hsqa{Thầy ơi em cần điểm để qua môn.}{Vậy em cần học để qua đời kiểm tra trước đã.}
\item \hsqa{Em xin làm bài khác bù được không?}{Bù được, nếu em chịu bù cả ý thức.}
\item \hsqa{Bạn kia điểm cao hơn em vô lý.}{Vô lý hay vì bạn đó ngồi học lúc em đang lướt điện thoại?}
\item \hsqa{Em cố rồi mà vẫn không tới.}{Cố là đáng quý, nhưng phải coi cố đúng hướng chưa.}
\item \hsqa{Điểm không nói lên tất cả mà thầy.}{Đúng, nhưng nó đang nói kha khá về bài này.}
\item \hsqa{Thầy ơi em buồn điểm này lắm.}{Buồn thì được, miễn là buồn xong chịu sửa.}
\item \hsqa{Em không dám đem điểm này về nhà.}{Vậy đem theo kế hoạch sửa, đỡ sợ hơn.}
\item \hsqa{Em xui thôi.}{Xui một ít, hổng kiến thức một ít, mình nói thật cho đủ.}
\item \hsqa{Thầy có thể bỏ qua lỗi này không?}{Bỏ qua trong đời được, trong bài chấm hơi khó.}
\item \hsqa{Em làm khác cách thầy dạy.}{Khác được, miễn đúng chứ không phải khác cho oách.}
\item \hsqa{Em không hợp cách chấm này.}{Bài thi cũng không cần hợp em mới phát huy tác dụng.}
\item \hsqa{Em lỡ tay xóa nhầm.}{Lỡ tay thì tiếc, nhưng điểm vẫn đi theo giấy.}
\item \hsqa{Thầy cho em cứu vớt chút đi.}{Cứu vớt bắt đầu từ em, không bắt đầu từ bút đỏ của thầy.}
\item \hsqa{Em làm được trong đầu mà quên ghi.}{Trong đầu em điểm đẹp, trên giấy em điểm thật.}
\item \hsqa{Em chỉ thiếu may mắn.}{May mắn không thay nổi phần ôn tập bị bỏ.}
\item \hsqa{Sao em học hoài vẫn không lên?}{Có khi em học hoài kiểu cũ nên kết quả cũng cũ.}
\item \hsqa{Em nghĩ mình bị chấm sót.}{Cứ kiểm tra kỹ, nhưng đừng lấy hy vọng thay cho dữ liệu.}
\end{enumerate}

\section{Từ câu 141 đến câu 160}
\begin{enumerate}[leftmargin=*,resume]
\item \hsqa{Thầy ơi em học tủ trúng ít quá.}{Tủ ít đau ít, tủ trật đau dài.}
\item \hsqa{Em không ngờ đề ra vậy.}{Đề cũng không ngờ em học vậy.}
\item \hsqa{Em được cộng chuyên cần không?}{Chuyên cần là nền, không phải phép màu.}
\item \hsqa{Sao em làm không tệ mà điểm tệ?}{Vì cảm giác không tệ và kết quả tệ đi cùng nhau hoài.}
\item \hsqa{Em thiếu đúng một ý.}{Ý đó đang giữ kha khá điểm.}
\item \hsqa{Em xin nợ điểm này được không?}{Điểm số không phải app trả góp.}
\item \hsqa{Thầy ơi có cách nào kéo điểm nhanh không?}{Có, bớt chống chế và học đều.}
\item \hsqa{Em đọc đề hiểu mà làm sai.}{Hiểu trong lòng khác với làm trên giấy.}
\item \hsqa{Em tiếc quá.}{Tiếc là tốt, miễn đừng chỉ tiếc bằng miệng.}
\item \hsqa{Điểm này không phản ánh em đâu.}{Không phản ánh hết em, nhưng phản ánh rõ bài thi này.}
\item \hsqa{Em không giỏi thi cử.}{Vậy càng phải giỏi chuẩn bị.}
\item \hsqa{Sao em làm giống bạn mà điểm thấp hơn?}{Giống bạn là em nghĩ vậy, bài chấm lại nghĩ khác.}
\item \hsqa{Em đáng lẽ cao hơn.}{Đáng lẽ là họ hàng của không xảy ra.}
\item \hsqa{Thầy châm chước cho em cuối kỳ đi.}{Cuối kỳ mới xin châm chước thì hơi muộn để dễ thương.}
\item \hsqa{Em tưởng câu đó không quan trọng.}{Câu em tưởng không quan trọng hay rất thích xuất hiện trên thang điểm.}
\item \hsqa{Em cần điểm đẹp để nhà vui.}{Vậy nhà vui bền thì em học bền trước.}
\item \hsqa{Em được 7 mà buồn hơn hồi được 5.}{Vì em bắt đầu thấy mình có thể hơn.}
\item \hsqa{Em cố mà không ai thấy.}{Điểm thấy, chỉ là chưa thấy đủ.}
\item \hsqa{Em có thể nhờ cảm tình không?}{Cảm tình tốt nhất là từ bài làm đúng.}
\item \hsqa{Điểm này làm em mất tinh thần.}{Mất một lúc được, mất luôn thì phí bài học.}
\end{enumerate}

\section{Từ câu 161 đến câu 180}
\begin{enumerate}[leftmargin=*,resume]
\item \hsqa{Thầy ơi em xin gỡ điểm danh dự.}{Danh dự gỡ bằng quá trình, không bằng câu xin.}
\item \hsqa{Em làm ẩu nên mất điểm oan.}{Không oan, bài em tố em rất rõ.}
\item \hsqa{Em không nghĩ điểm lại thấp vậy.}{Nhiều lỗ hổng cũng nghĩ em không nhận ra nó.}
\item \hsqa{Thầy cho em câu động viên đi.}{Động viên đây: lần sau học sớm.}
\item \hsqa{Em hứa lần sau sẽ khác.}{Tốt, miễn lần sau đừng chỉ khác ở lời hứa.}
\item \hsqa{Em chỉ cần đủ qua là được.}{Qua được thì ổn, nhưng sống kiểu vừa đủ mãi hơi nguy hiểm.}
\item \hsqa{Thầy đừng công bố điểm to quá.}{Điểm nhỏ mà bài làm em to chuyện quá.}
\item \hsqa{Em muốn quên điểm này.}{Quên cảm giác được, đừng quên nguyên nhân.}
\item \hsqa{Em học lệch nên bị vậy.}{Lệch thì điểm kéo thẳng lại thôi.}
\item \hsqa{Thầy cho em cơ hội cuối.}{Cuối mấy lần rồi em?}
\item \hsqa{Em cần động lực để kéo điểm.}{Động lực lớn nhất đang nằm trên bảng điểm đó.}
\item \hsqa{Em sợ xem bài quá.}{Sợ mà vẫn xem mới lớn.}
\item \hsqa{Em không dám đối diện.}{Không đối diện thì lỗi vẫn ở đó chờ em.}
\item \hsqa{Em gắng cả tuần cho bài này.}{Gắng đáng khen, nhưng sai vẫn cần sửa.}
\item \hsqa{Em thiếu cẩn thận thôi.}{Thiếu cẩn thận trong thi cử là thiếu đúng chỗ đau nhất.}
\item \hsqa{Thầy ơi em được nâng tròn không?}{Điểm chứ có phải giá rau đâu mà nâng tròn.}
\item \hsqa{Em tưởng mình làm ổn.}{Tưởng ổn là bệnh khá phổ biến ở học sinh.}
\item \hsqa{Sao em cứ kém ở phần này?}{Vì em cứ tránh nhìn đúng phần này.}
\item \hsqa{Em cần điểm đẹp để tự tin.}{Tự tin bền phải đi cùng thực lực.}
\item \hsqa{Thầy ơi bài em còn cứu được không?}{Cứu được, nếu chủ bài chịu cứu mình trước.}
\end{enumerate}

\section{Từ câu 181 đến câu 200}
\begin{enumerate}[leftmargin=*,resume]
\item \hsqa{Em thấy mình không có duyên với điểm cao.}{Duyên ít khi thắng được nề nếp.}
\item \hsqa{Em học bài mà vô phòng thi là trắng.}{Vậy em cần luyện tình huống, không chỉ luyện trí nhớ.}
\item \hsqa{Em có thể xin kiểm tra lại không?}{Được, miễn em kiểm tra lại luôn cách học của em.}
\item \hsqa{Điểm này làm em quê quá.}{Quê chút cho nhớ lâu cũng đáng.}
\item \hsqa{Em thiếu đúng vài chữ.}{Vài chữ nhiều khi là vài điểm.}
\item \hsqa{Thầy ơi sao em mãi không bứt lên được?}{Vì em còn đứng trên cái nền cũ.}
\item \hsqa{Em học dàn trải quá.}{Ừ, nên điểm em cũng dàn đều ở mức lưng chừng.}
\item \hsqa{Em không ngờ bài dễ vậy mà vẫn thấp.}{Đó mới là điều đáng suy nghĩ nhất.}
\item \hsqa{Em muốn hỏi cách lên điểm nhanh.}{Nhanh nhất là đừng chờ nước tới chân mới học.}
\item \hsqa{Em hay bị mất điểm lặt vặt.}{Lặt vặt góp lại thường đấm khá đau.}
\item \hsqa{Em sợ nhà la vì điểm này.}{Sợ thì biến nó thành kế hoạch sửa, đừng biến thành bí mật.}
\item \hsqa{Em có thể được cứu xét đặc biệt không?}{Đặc biệt nhất là em chịu học khác đi.}
\item \hsqa{Em muốn bài được chấm rộng lượng hơn.}{Rộng lượng quá thì em lại chậm lớn.}
\item \hsqa{Thầy ơi em không đáng thấp vậy đâu.}{Thầy cũng mong thế, nhưng bài chưa đồng ý.}
\item \hsqa{Điểm này làm em nản học.}{Điểm này nên làm em tỉnh học.}
\item \hsqa{Em được bao nhiêu thì vừa ý thầy?}{Thầy vừa ý nhất khi em đúng phần đáng đúng.}
\item \hsqa{Em có thể làm thêm bài lấy điểm không?}{Có, miễn em không xem đó là vé né quá trình.}
\item \hsqa{Em thấy áp lực vì điểm quá.}{Áp lực vừa đủ thì kéo em đi, còn né tránh thì kéo em xuống.}
\item \hsqa{Em có hi vọng cứu học kỳ này không?}{Hi vọng có, miễn em bỏ thói quen trễ nải của học kỳ trước.}
\item \hsqa{Thầy ơi chốt lại em cần làm gì để điểm lên?}{Học đều, làm thật, sửa kỹ, bớt xin.}
\end{enumerate}
